%%
%% fall-lecture02.tex
%% 
%% Made by Alex Nelson <pqnelson@gmail.com>
%% Login   <alex@lisp>
%% 
%% Started on  2026-10-13T11:39:11-0700
%% Last update 2026-10-13T11:39:11-0700
%% 

\lecture{}

% Section 2: Formal construction and fundamental examples
% Subsection 2.1: Vector bundles

\subsection{Vector bundles}

\begin{definition}[Real vector bundle]
Let $n\in\NN$ be a natural number.
A \define{Real Vector Bundle} of rank $n$, $\xi$, consists of
\begin{enumerate}
\item a topological space $E=E(\xi)$ called the ``total space'',
\item a ``base space'' $B$ and a continuous projection map $\pi\colon E\to B$,
\item for each $b\in B$, each fiber $F_{b}=\pi^{-1}\{b\}$ has a real
  vector space structure
\end{enumerate}
such that
\begin{enumerate}
\item Local triviality: for all $b\in B$ there exists an open subset
  $U\subset B$ such that $b\in U$ and there exists a continuous
  $h\colon U\times\RR^{n}\to\pi^{-1}(U)$ such that $h$ is a
  homeomorphism and for all $b'\in U$ we have $h(b',-)\colon\RR^{n}\to F_{b}$
  is a vector space isomorphism. We call such an $h$ a \define{Local Trivialization}
  of $\xi$.
\end{enumerate}
The pair $(U,h)$ is a \define{Local Coordinate System} for $\xi$
around $b$.
\end{definition}

\begin{remark}
Replace $\RR^{n}$ by $\CC^{n}$, we get the definition of a complex
vector bundle.
\end{remark}

\begin{remark}
When $n=1$, rank-$n$ bundles are just called \define{Line Bundles}.
\end{remark}

\begin{node}
We also have a notion of $C^{\infty}$-vector bundle (pronounced
``smooth vector bundle''). Just replace ``continuous'' by ``smooth',
and ``homeomorphism'' by ``diffeomorphism'', and require $B$ to be a
smooth manifold.
\end{node}

\begin{example}[Trivial bundle]
The trivial bundle $B\times\RR^{n}=E$, the projection map $\pi(b,v)=b$.
The local triviality condition si satisfied by taking $U:=B$,
$h:=\id$. The vector space structure on each fiber looks like:
\begin{equation}
c_{1}(b,v_{1})+c_{2}(b,v_{2})=(b,c_{1}v_{1}+c_{2}v_{2})
\end{equation}
for any $c_{1},c_{2}\in\RR$.
\end{example}

\begin{example}[Tangent bundle of smooth manifold]
Let $M$ be a smooth $n$-manifold. Then
\begin{equation}
\T M=\{(x,v)\mid x\in M,v\in \T_{x}M\}.
\end{equation}
Then the projection map is $\pi(x,v)=x$.

How to check it is locally trivial/ We need a smooth chart of $\T M$.
First, choose $(U,\varphi)$ to be a chart of $M$. Then
$(\pi^{-1}(U),\Phi)$ where for any tangent vector
$v=(x,\sum_{j}\xi^{j}\partial_{j}|_{x})$ we have
\begin{equation}
\Phi(v)=(\varphi(\pi(v)),\vec{\xi})
\end{equation}
where $\vec{\xi}=(\xi^{1},\dots,\xi^{n})$ are the coefficients. The
local trivialization $(U,h)$ where $h\colon U\times\RR^{n}\to\pi^{-1}(U)$.
Well, $h(x,\xi)=v$. We need to prove it is a diffeomorphism (but
that's obvious).
\end{example}

\begin{definition}
Let $\xi$ and $\eta$ be vector bundles over the same base space $B$
(both $\xi$ and $\eta$ real vector bundles, or more generally
both with fibers are vector spaces over the same field).
We say $\xi$ is \define{Isomorphic}
to $\eta$, denoted $\xi\iso\eta$, if there exists a homeomorphism
$f\colon E(\xi)\to E(\eta)$ between the total spaces such that each
vector space $F_{b}(\xi)$ is mapped isomorphically onto the
corresponding vector space $F_{b}(\eta)$.

\textbf{Note:} both $\xi$ and $\eta$ are vector bundles over
\emph{the same} base space $B$. This is how Milnor and
Stasheff~\cite[Ch~2]{milnor1974characteristic} define it. But I am not
sure how I feel about it: does this correspond to the notion of an
isomorphism in the category of vector bundles? (That is to say, using
Definition~\ref{defn:bundle-map} of bundle maps, and a bundle
isomorphism is a bundle map with an inverse bundle map, we could have
homeomorphic base spaces $B(\xi)\iso B(\eta)$ which are not
necessarily identical---is this a ``bad thing''?)
\end{definition}

\begin{definition}
Let $M$ be a smooth $n$-manifold.
We say $M$ is \define{Parallelizable} if $\T M$
is trivializable (i.e., globally trivial $\T M=M\times\RR^{n}$).
\end{definition}

\begin{example}
Recall, for any $n$-manifold, there is the Whitney embedding
$M\into\RR^{n}$ where $N>n$ is some suitably bigger number. THen we
have the \define{Normal Bundle}
\begin{equation}
NM = \{(x,v)\mid x\in M,v\in(\T_{x}M)^{\perp}\}.
\end{equation}
\end{example}

\begin{example}[Canonical line bundle of $\RP^{n}$]
Recall $\RP^{n}$ is the set of lines passing through the origin in
$\RR^{n}$. We have
\begin{align*}
q\colon & \sphere{n}\to\RP^{n}\\
& x\mapsto[x]=\{\pm x\}.
\end{align*}
Then the canonical line bundle
\begin{equation}
E(\gamma^{1}_{n}):=\{([x],v)\mid [x]\in\RP^{n},v\in\RR^{n+1},\exists\lambda\in\RR\setminus\{0\}\ldotp
v=\lambda x\}.
\end{equation}
We will still need to check it has a local triviality. First, choose
some open set for any $[x]\in\RP^{n}$, we have $\{\pm x\}\subset\sphere{n}$,
then choose some small open neighborhood $U\subset\RP^{n}$ such that
$q^{-1}U$ is some suitable small neighborhood . Then, letting $U_{1}=q(U)$,
\begin{align*}
h\colon & U_{1}\times\RR\to\pi^{-1}(U_{1})\\
& ([x],t)\mapsto([x],tx)
\end{align*}
where $x\in U_{+}$ is in one of the connected components of $U$.
\end{example}

\begin{definition}[Cross section]
Let $\xi$ be a bundle $E(\xi)$ with base $B$.
A \define{Cross Section} of $\xi$ is a continuous map $s\colon B\to E(\xi)$
such that $\pi\circ s=\id$.
\end{definition}

\begin{remark}
If $B$ is a smooth manifold, and if $E=\T B$ is the tangent bundle,
then a cross section is just a vector field.
\end{remark}

\begin{definition}
A cross section is \define{Nowhere Zero} if for all $b\in B$ we have $s(b)\neq0$.
\end{definition}

\begin{theorem}
A line bundle $\xi$ is trivial if and only if it has a nowhere zero section.
\end{theorem}

\begin{proof}
\forwardproof\ We see $E(\xi)\iso B\times\RR$ and then we just take $s(b)=(b,1)$
which is the desired nowhere zero section.

\backwardproof\ Assume we are given a nowhere zero cross section
$s\colon B\to E(\xi)$ where for each $b\in B$ we have $s(b)=(b,e(b))$
where $e(b)\neq0$. For each $v\in\pi^{-1}(b)$, we can write
$v=\lambda^{v}_{b}\cdot e(b)$ [where $\lambda^{v}_{b}\in\RR$ is just
the nonzero coefficient],
then $E(\xi)\to B\times\RR^{n}$ maps $v\mapsto(b,\lambda^{v}_{b})$.
\end{proof}

\begin{example}
The tangent bundle $\T\sphere{1}$ is a trivial bundle; i.e.,
$\sphere{1}$ is trivializable. Then
\begin{equation}
s\bigl((x_{1},x_{2})\bigr)=\bigl((x_{1},x_{2}),(-x_{2},x_{1})\bigr)
\end{equation}
is the desired nowhere zero section.
\end{example}

\begin{non-example}
The canonical line bundle $\gamma^{1}_{n}$ is nontrivial for every $n\geq1$.
Assume for contradiction $s\colon\RP^{n}\to E(\gamma^{1}_{n})$ is the
nowhere vanishing section. Then
\begin{align*}
\sphere{n} & \xrightarrow{q}\RP^{n}\xrightarrow{s} E(\gamma^{1}_{n})\\
x&\mapsto [x]\mapsto([x],t(x)x),
\end{align*}
where $t(x)$ is continuous and non-vanishing. But $\sphere{n}$ is
connected and
\begin{subequations}
\begin{equation}
x\mapsto([x],t(x)x),
\end{equation}
and
\begin{equation}
-x\mapsto([x],t(-x)(-x)),
\end{equation}
\end{subequations}
are mapped to the same point. This means
\begin{equation}
t(x)=-t(-x),
\end{equation}
so $t$ is an odd function. Then by the intermediate value theorem,
there must be an $x_{0}\in\sphere{n}$ such that $t(x_{0})=0$,
contraditing the non-vanishing assumption.
\end{non-example}

\begin{definition}
A set of cross sections $\{s_{1},\dots,s_{n}\}$ is \define{Nowhere Dependent}
if $s_{1}(b)$, \dots, $s_{n}(b)$ are linearly independent for each
$b\in B$. (This is a generalization of nowhere vanishing.)
\end{definition}

\begin{theorem}
An $\RR^{n}$-vector bundle is trivial if and only if there exists
sections $s_{1}$, \dots, $s_{n}$ which is nowhere dependent.
\end{theorem}

(The proof is similar to the previous theorem's proof.)

\begin{example}
The tangent bundle $\T\sphere{3}$ is trivial. Writing $x=(x_{1},x_{2},x_{3},x_{4})\in\sphere{3}\subset\RR^{4}$,
then
\begin{subequations}
  \begin{align}
s_{1}(x) &= (x, (-x_{2},x_{1},-x_{4},x_{3}))\\
s_{2}(x) &= (x, (-x_{3},x_{2},+x_{1},-x_{4}))\\
s_{3}(x) &= (x, (-x_{4},x_{3},-x_{2},x_{3}))
  \end{align}
\end{subequations}
and you cancheck they are linearly independent.
\end{example}

\begin{theorem}
The only parallelizable spheres are $\sphere{0}$, $\sphere{1}$,
$\sphere{3}$, $\sphere{}$. In particular, $\T\sphere{2}$ is not trivial.
\end{theorem}

(This is a corollary of the hairy ball theorem.)